\section{Architecture}

\subsection{Component Map}

The codebase consists of seven source modules, each with a single well-defined responsibility:

\begin{center}
\begin{tabular}{>{\ttfamily}l l}
\toprule
\textbf{Module} & \textbf{Responsibility} \\
\midrule
args (.hpp / .cpp)       & CLI argument parsing; populates \texttt{ParsedArgs} \\
bitio (.hpp only)        & Bit-level read/write abstraction over byte streams \\
lzss (.hpp / .cpp)       & LZSS sliding-window encoder and decoder \\
model (.hpp / .cpp)      & Forward and inverse first-order delta transform \\
scanner (.hpp / .cpp)    & Image-to-stream reordering and block decomposition \\
codec (.hpp / .cpp)      & Top-level compress/decompress orchestration + file format \\
main (.cpp only)         & Entry point; opens streams, dispatches to codec \\
\bottomrule
\end{tabular}
\end{center}

\subsection{Compression Data Flow}

The following steps occur during a compression call:

\begin{enumerate}
  \item \texttt{main.cpp} calls \texttt{parse\_arguments(argc, argv)} $\to$ \texttt{ParsedArgs}.
  \item Opens \texttt{args.infile} as a binary \texttt{std::ifstream} and \texttt{args.outfile}
        as a binary \texttt{std::ofstream}.
  \item Calls \texttt{compress(in, out, args.width, args.use\_model, args.adaptive\_scan)}.
  \item \texttt{codec.cpp}: reads all bytes from \texttt{in} into \texttt{vector\textless uint8\_t\textgreater{} pixels}.
        Validates that \texttt{pixels.size() \% width == 0}.
  \item Writes the LZKO file header (magic, width, height, flags, block\_size\_log2, original size).
  \item \textbf{Static path} (no \texttt{-a}):
    \begin{enumerate}
      \item \texttt{scan\_horizontal(pixels, w, h)} $\to$ 1-D byte stream (trivial copy).
      \item If \texttt{-m}: \texttt{forward\_model(data)} transforms in-place.
      \item \texttt{lzss\_encode(data, BitWriter(out))}; then \texttt{bw.flush()}.
    \end{enumerate}
  \item \textbf{Adaptive path} (\texttt{-a}), two passes:
    \begin{enumerate}
      \item \textbf{Pass 1 --- metadata}: \texttt{split\_into\_blocks(pixels, w, h, 16)} $\to$ \texttt{blocks}.
            For each block: compute SAD in both directions; choose lower-SAD direction as
            \texttt{scan\_dir}; scan; optionally \texttt{forward\_model}; \texttt{lzss\_encode\_to\_buffer}
            to get compressed size; set \texttt{meta.compressed} = 1 if compressed \textless{} raw.
      \item Write \texttt{num\_blocks} and packed 2-bit-per-block metadata to output.
      \item \textbf{Pass 2 --- data}: repeat scan + model + encode for each block;
            write \texttt{uint32\_t stored\_size} followed by the block bytes.
    \end{enumerate}
\end{enumerate}

\subsection{Decompression Data Flow}

The following steps occur during a decompression call:

\begin{enumerate}
  \item \texttt{main.cpp} calls \texttt{decompress(in, out)}.
  \item \texttt{codec.cpp}: reads and validates magic bytes; returns error code 1 on mismatch.
  \item Reads header fields: \texttt{width}, \texttt{height}, \texttt{flags}, \texttt{block\_size\_log2},
        \texttt{original\_size}. Reconstructs \texttt{use\_model} and \texttt{adaptive} from \texttt{flags}.
  \item \textbf{Static path}:
    \begin{enumerate}
      \item \texttt{lzss\_decode(BitReader(in), data, original\_size)}.
      \item If \texttt{use\_model}: \texttt{inverse\_model(data)}.
      \item Write \texttt{data} to \texttt{out}.
    \end{enumerate}
  \item \textbf{Adaptive path}:
    \begin{enumerate}
      \item Read \texttt{num\_blocks} and packed block metadata; unpack into \texttt{scan\_dirs[]}
            and \texttt{comp\_flags[]}.
      \item Allocate full image buffer \texttt{vector\textless uint8\_t\textgreater{} image(width * height)}.
      \item For each block index \texttt{idx}: read \texttt{stored\_size} + raw bytes;
            if \texttt{comp\_flags[idx]}: wrap in \texttt{istringstream} + \texttt{BitReader} +
            \texttt{lzss\_decode}; else use raw bytes directly.
      \item If \texttt{use\_model}: \texttt{inverse\_model(block\_data)}.
      \item \texttt{unscan\_horizontal} or \texttt{unscan\_vertical} depending on \texttt{scan\_dirs[idx]}.
      \item Copy block pixels back into correct position in \texttt{image} buffer.
      \item After all blocks: write \texttt{image} to \texttt{out}.
    \end{enumerate}
\end{enumerate}

\subsection{Dependency Graph}

\begin{center}
\begin{tabular}{>{\ttfamily}l >{\ttfamily}l}
\toprule
\textbf{Module} & \textbf{Depends on} \\
\midrule
args     & (none) \\
bitio    & (none) \\
model    & (none) \\
scanner  & (none) \\
lzss     & bitio \\
codec    & bitio, lzss, model, scanner \\
main     & args, codec \\
\bottomrule
\end{tabular}
\end{center}

\subsection{File Format}

Every compressed file begins with the LZKO header:

\begin{center}
\begin{tabular}{r r l}
\toprule
\textbf{Offset} & \textbf{Size} & \textbf{Field} \\
\midrule
0  & 4 B & Magic bytes: \texttt{0x4C 0x5A 0x4B 0x4F} (``LZKO'') \\
4  & 2 B & \texttt{uint16\_t} image width (little-endian) \\
6  & 2 B & \texttt{uint16\_t} image height (little-endian) \\
8  & 1 B & Flags: bit 0 = use\_model, bit 1 = adaptive \\
9  & 1 B & \texttt{block\_size\_log2} (= 4 $\to$ $2^4 = 16$ pixels) \\
10 & 4 B & \texttt{uint32\_t} original pixel count (little-endian) \\
14 & \ldots & Compressed data (static) \emph{or} block section (adaptive) \\
\bottomrule
\end{tabular}
\end{center}

In adaptive mode, the data section is structured as:
\begin{itemize}
  \item \texttt{uint32\_t num\_blocks} --- total number of 16$\times$16 blocks.
  \item Bit-packed metadata: 2 bits per block, MSB-first, padded to a byte boundary.
        Bit~0 of each pair = \texttt{scan\_dir}; bit~1 = \texttt{compressed}.
  \item Per-block data: \texttt{uint32\_t stored\_size} followed by \texttt{stored\_size} bytes
        (LZSS-compressed or raw).
\end{itemize}
